%% 03-results.tex — Three bullet results

\section{What the Evidence Shows}

Applied to all fifty states across the 2000, 2010, and 2020 Census cycles,
the algorithm produces the following results:

\bigskip

\noindent
\begin{tabular}{@{}>{\bfseries\large}p{0.08\textwidth} p{0.86\textwidth}@{}}
+22\% &
More compact than enacted 2020 maps, measured by the Polsby-Popper
ratio of district area to perimeter squared, averaged across all 435 districts.
Compact districts follow rivers and county lines rather than
twisting through neighborhoods to sort voters.
\textit{(Paper B.2~\cite{dellaLibera2026edgeWeighted})}\\[10pt]

7D / 7R &
North Carolina---the state that has generated more federal gerrymandering
litigation than any other in the past decade---receives a
seven-Democratic, seven-Republican delegation under the algorithm,
closely matching the state's evenly divided electorate.
No partisan data was used.
\textit{(Paper B.11~\cite{dellaLibera2026ncRegions})}\\[10pt]

15--20 min &
Any citizen with a laptop and an internet connection can download
the open-source \texttt{redist} tool, fetch public Census data,
run the Vermont verification fixture, and confirm the result themselves
in 15--20 minutes end-to-end including data download.
Full fifty-state replication takes approximately 18 minutes for one
census year, or two to four hours for all three census years.
\textit{(Replication Package A.4)}\\
\end{tabular}

\bigskip

\noindent Additional findings:
the algorithm produces \textbf{more majority-minority districts}
than enacted plans in states above the 42\% minority-population threshold
identified in Paper D.1~\cite{dellaLibera2026vraCompliance},
and reduces the partisan efficiency gap by \textbf{62\%} compared to
enacted plans~\cite{dellaLibera2026efficiencyGap}---without any
partisan fairness criterion in the algorithm specification.\footnote{%
The algorithm produces near-zero efficiency gaps as a byproduct of geometric
neutrality (mean $|EG|=0.04$ vs.\ $0.08$ for enacted maps).
In some geographically sorted states a small systematic gap persists due to
urban concentration (Paper C.5).}
